\section{Phương trình quy về phương trình bậc hai}

\subsection{Đề 1}

\Opensolutionfile{ans}[ans/ans-K10_C6-B4-D1]
\caulc

\begin{ex}%[0D7H3-1]
	Tổng tất cả các nghiệm của phương trình $\sqrt{x^2 + 2x - 3} = \sqrt{15-5x}$ là
	\choice
	{$7$}
	{\True $-7$}
	{$6$}
	{$4$}
	\loigiai{
	Bình phương hai vế của phương trình ta được 
	\[
		x^2 + 2x -3 = 15-5x \Rightarrow x^2  + 7x-18 = 0\Rightarrow \hoac{&x=2\\ &x=-9.}
	\]	
	Thử lại ta thấy $x=2$ và $x=-9$ đều thoả mãn nên tổng các nghiệm của phương trình là 
	\[
		2+(-9)= -7.
	\]
}
\end{ex}

\begin{ex}%[0D7H3-2]
	Tập nghiệm $S$ của phương trình $\sqrt{x^2-4} = x-2 $ là 
	\choice
	{$S = \{0\}$}
	{\True $S = \{2\}$}
	{$S = \{0;2\}$}
	{$S = \emptyset$}
	\loigiai{
		Ta có phương trình tương đương với
	\[
		\Leftrightarrow \heva{&x-2\ge 0\\ &x^2-4 = (x-2)^2} \Leftrightarrow \heva{&x\ge 2\\ &4x-8 = 0} \Leftrightarrow x=2.
	\]	
	Vậy tập nghiệm của phương trình là $S = \{2\}$.
}
\end{ex}

\begin{ex}%[0D7N3-2]
	Cho phương trình $\sqrt{3x+4} = x\quad (1)$. Mệnh đề nào sau đây đúng?
	\choice
	{$(1)\Leftrightarrow 3x+4=x^2$}
	{$(1)\Leftrightarrow 3x+4=x$}
	{$(1)\Leftrightarrow \heva{&3x+4\ge 0\\ &3x+4=x^2}$}
	{\True $(1)\Leftrightarrow \heva{&x\ge 0\\ &3x+4=x^2}$}
	\loigiai{
	Ta có $(1)\Leftrightarrow \heva{&x\ge 0\\ &3x+4=x^2}$	
}
\end{ex}

\begin{ex}%[0D7H3-2]
	Tập nghiệm $S$ của phương trình $\sqrt{2x-3} = x-3 $ là 
	\choice
	{\True $S = \{6\}$}
	{$S = \{2\}$}
	{$S = \{2;6\}$}
	{$S = \emptyset$}
	\loigiai{
		Ta có phương trình tương đương với
		\[
		\Leftrightarrow \heva{&x-3\ge 0\\ &2x-3 = (x-3)^2} \Leftrightarrow \heva{&x\ge 3\\ &x^2 - 8x + 12 = 0} \Leftrightarrow \heva{&x\ge 3\\ &\hoac{&x=2\\ &x=6}} \Leftrightarrow x=6.
		\]	
		Vậy tập nghiệm của phương trình là $S = \{6\}$.
	}
\end{ex}

\begin{ex}%[0D7H3-2]
	Số nghiệm của phương trình $\sqrt{3x^2 - 9x + 7} = x - 2$ là
	\choice
	{$3$}
	{$1$}
	{$2$}
	{\True $0$}
	\loigiai{
	Ta có phương trình tương đương với
	\[
		\heva{&x-2\ge 0\\ &3x^2 - 9x + 7 = (x-2)^2}
		\Leftrightarrow \heva{&x\ge 2\\ &2x^2 - 5x+3=0}
		\Leftrightarrow \heva{&x\ge 2\\ &\hoac{&x=1\\ &x=\dfrac{3}{2}}}
		\Leftrightarrow x\in \emptyset.
	\]	
	Vậy phương trình vô nghiệm.
}
\end{ex}

\begin{ex}%[0D7H3-2]
	Tập nghiệm của phương trình $\sqrt{5x+6} = x-6$ là 
	\choice
	{$S=\{2\}$}
	{$S = \{6\}$}
	{$S = \{2;15\}$}
	{\True $S = \{15\}$}
	\loigiai{
	Ta có phương trình tương đương với
	\[
		\heva{&x-6\ge 0\\ &5x+6 = (x-6)^2}
		\Leftrightarrow \heva{&x\ge 6\\ &x^2 - 17x + 30 = 0}
		\Leftrightarrow \heva{&x\ge 6\\ &\hoac{&x=2\\ &x=15}}
		\Leftrightarrow x=15.
	\]	
	Vậy tập nghiệm của phương trình là $S = \{ 15 \}$.
}
\end{ex}

\begin{ex}%[0D7H3-1]
	Số nghiệm của phương trình $\sqrt{x^2 + x - 42} = \sqrt{2x - 30}$ là 
	\choice
	{\True $0$}
	{$1$}
	{$2$}
	{vô số}
	\loigiai{
	Bình phương hai vế của phương trình ta được 
	\[
		x^2 + x - 42=2x - 30 \Leftrightarrow x^2-x-12=0 \Leftrightarrow 
		\hoac{&x=4\\&x=-3.} 
	\]
	Thay hai giá trị $x=-3$ và $x=4$ vào phương trình ban đầu ta thấy không có giá trị nào thõa mãn phương trình.\\
	Vậy phương trình vô nghiệm.	
}
\end{ex}

\begin{ex}%[0D7H3-1]
	Tập nghiệm của phương trình $2\sqrt{x^2 - x - 1} = \sqrt{x^2 + 2x + 5}$ là
	\choice
	{$S = \{1\}$}
	{$S = \{3\}$}
	{\True $S=\{1;3\}$}
	{$S = \emptyset$}
	\loigiai{
	$2\sqrt{x^2 - x - 1} = \sqrt{x^2 + 2x + 5}.\quad (1)$\\
	Bình phương hai vế của phương trình $(1)$ ta được phương trình $4(x^2 - x - 1) = x^2 + 2x + 5$\\
	Ta có $4(x^2 - x - 1) = x^2 + 2x + 5 \Leftrightarrow 3x^2-6x-9=0 \Leftrightarrow \hoac{&	x=-1\\&	x=3.}$\\
	Thay hai giá trị $x=-1$ và $x=3$ vào phương trình $(1)$ ta thấy cả hai giá trị đều thõa mãn. Vậy phương trình có hai nghiệm là $x=-1$ và $x=3$.	
}
\end{ex}

\begin{ex}%[0D7H3-2]
	Số nghiệm của phương trình $\sqrt{3x^2 - 9x + 1} = |x-2|$ là
	\choice
	{$0$}
	{$1$}
	{\True $2$}
	{$3$}
	\loigiai{
		Ta có phương trình tương đương với 
		\[
		3x^2 - 9x + 1  = (x-2)^2 \Leftrightarrow 2x^2 - x - 21 = 0 \Leftrightarrow \hoac{&x=-3 \\ & x = \dfrac{7}{2}.}
		\]
		Vậy phương trình có $2$ nghiệm.
	}
\end{ex}

\begin{ex}%[0D7H3-1]
	Nghiệm của phương trình $\sqrt{x^2 + 3x - 2} = \sqrt{x + 1}$ là
	\choice
	{\True $x = 1$}
	{$x= -3$}
	{$\hoac{& x = -3 \\ & x = 1}$}	  
	{$ x \in \emptyset $}
	\loigiai{
		Điều kiện: $\heva{&x^2 + 3x - 2 \geq 0\\&x + 1 \geq 0}$.\\
		Ta có $\sqrt{x^2 + 3x - 2} = \sqrt{x + 1} \Leftrightarrow x^2 + 3x - 2 = x + 1 \Leftrightarrow x^2 + 2x - 3 = 0 \Leftrightarrow \hoac{&x = 1\\&x = -3 \;\;\text{(loại)}}$.\\
		Vậy phương trình đã cho có nghiệm là $x = 1$.
	}
\end{ex}

\begin{ex}%[0D7H3-1]
	Nghiệm của phương trình $\sqrt{x^2 - 5x + 4} = \sqrt{-2x^2 - 3x + 12}$ là
	\choice
	{\True $ x = 2$}	
	{$ x = -\dfrac{4}{3} $}
	{$\hoac{& x = 2\\& x = -\dfrac{4}{3}}$} 
	{$ x \in \emptyset$}
	\loigiai{
		Điều kiện: $\heva{&x^2 - 5x + 4 \geq 0\\&-2x^2 - 3x + 12 \geq 0}$.\\
		Ta có
		\allowdisplaybreaks
		\begin{eqnarray*}
			\sqrt{x^2 - 5x + 4} = \sqrt{-2x^2 - 3x + 12}
			&\Leftrightarrow&x^2 - 5x + 4 = -2x^2 - 3x + 12\\
			&\Leftrightarrow&3x^2 - 2x - 8 = 0\\
			&\Leftrightarrow&\hoac{&x = 2 \,\text{ (loại)}\\&x = -\dfrac{4}{3}\, \text{ (thoả mãn)}.}
		\end{eqnarray*}
		Vậy phương trình đã cho có nghiệm là $x = 2$.
	}
\end{ex}



\begin{ex}%[0D7H3-2]
	Nghiệm của phương trình $\sqrt{4x^2+2x+10}=3x+1$ là
	\choice
	{\True $x=1$}
	{ $x=\dfrac{-9}{5}$}
	{$\hoac{&x=\dfrac{-9}{5}\\&x=1}$}
	{$x\in\emptyset$}
	\loigiai{Bình phương hai vế của phương trình đã cho, ta được
		\begin{eqnarray*}
			4x^2+2x+10=(3x+1)^2 \Rightarrow -5x^2-4x+9=0\Rightarrow \hoac{&x=1\\ &x=-\dfrac{9}{5}.}
		\end{eqnarray*}
		Thay lần lượt các giá trị trên vào phương trình đã cho, ta thấy $x=1$ thỏa mãn.\\
		Vậy nghiệm của phương trình đã cho là $x=1$.
	}
\end{ex}

\Closesolutionfile{ans}
\indapan{6}{ans/ans-K10_C6-B4-D1}
\cauds
\Opensolutionfile{ans}[ans/ans-K10_C6-B4-D1-DS]

\begin{ex}%[0D7V3-2]
	Cho phương trình $\sqrt{2x^2 + x+ 3} = -(x+5)\quad (1)$. Xác định tính Đúng/Sai của các mệnh đề sau.
	\choiceTF[t]
	{\True Bình phương hai vế của phương trình $(1)$ ta được phương trình $x^2 - 9x - 22 = 0$}
	{Phương trình $(1)$ cùng tập nghiệm với phương trình $x^2 - 9x - 22 = 0$}
	{$x=11$ và $x=-2$ là hai nghiệm của phương trình $(1)$}
	{\True Phương trình $(1)$ vô nghiệm}
	\loigiai{
		Bình phương hai vế của phương trình $(1)$ ta được 
		\[
			2x^2 + x + 3  = (x+5)^2 \Leftrightarrow x^2 - 9x - 22 = 0 \Leftrightarrow \hoac{&x=11\\ & x = -2.}
		\]
		Thử lại ta thấy cả $x=11$ và $x=-2$ đều không thoả mãn. Vậy
		\begin{itemchoice}
			\itemch Mệnh đề đúng. 
			\itemch Phương trình $(1)$ vô nghiệm mà tập nghiệm của phương trình $x^2 - 9x - 22 = 0$ là $S  = \{ -2;11 \}$.\\
			Mệnh đề sai.
			\itemch Mệnh đề sai. 
			\itemch Mệnh đề đúng. 
		\end{itemchoice}	
	}
\end{ex}

\begin{ex}%[0D7V3-1]
	Cho phương trình $\sqrt{x^2 + 2x + 4} = \sqrt{2-x}$. Xác định tính Đúng/Sai của các mệnh đề sau.
	\choiceTF[t]
	{\True Điều kiện xác định của phương trình $(1)$ là $x\le 2$}
	{Phương trình $(1)$ có hai nghiệm phân biệt trái dấu}
	{\True Phương trình $(1)$ có các nghiệm đều là số nguyên.}
	{Tổng các nghiệm của phương trình $(1)$ là $1$}
	\loigiai{
		Ta có 
		\[
			(1)\Leftrightarrow \heva{&x^2 + 2x + 4 \ge 0\\ &2-x\ge 0\\ & x^2 + 2x + 4  = 0}
			\Leftrightarrow \heva{& x\le 2\\ &x^2 + 3x + 2 = 0}
			\Leftrightarrow \hoac{&x=-1\\ &x=-2.}
		\]
		\begin{itemchoice}
			\itemch Điều kiện xác định của phương trình là $
				\heva{&x^2 + 2x + 4 \ge 0\\ &2-x\ge 0}
				\Leftrightarrow x\le 2.
			$\\
			Mệnh đề đúng. 
			\itemch Phương trình $(1)$ có hai nghiệm phân biệt đều âm.\\
			Mệnh đề sai.
			\itemch $-1$ và $-2$ đều là các số nguyên.\\
			Mệnh đề đúng.
			\itemch $-1+(-2) = -3$.\\
			Mệnh đề sai.
		\end{itemchoice}	
	}
\end{ex}

\begin{ex}%[0D7V3-1]
	Cho phương trình $(x+1)\left( \sqrt{x+4} - \sqrt{-x^2 + 4x + 14} \right) = 0 \quad (1)$. Xác định tính Đúng/Sai của các mệnh đề sau.
	\choiceTF[t]
	{Phương trình $(1)$ không có nghiệm bé thua $4$}
	{\True Phương trình $(1)$ có $3$ nghiệm phân biệt}
	{\True Các nghiệm của phương trình $(1)$ đều là số nguyên}
	{Tổng bình phương các nghiệm của phương trình $(1)$ bằng $29$}
	\loigiai{
		Ta có 
		\[
			(1) \Rightarrow \hoac{&x+1 = 0 \\ & \sqrt{x+4} - \sqrt{-x^2 + 4x + 14} = 0}
			\Rightarrow \hoac{&x=-1\\ &x+4  = -x^2 + 4x + 14}
			\Leftrightarrow \hoac{&x=-1\\ &x^2 - 3x - 10 = 0.}  
		\]
		Từ đó ta giải được các nghiệm $x=-1$, $x=-2$, $x=5$.\\
		Thử lại ta thấy $x=-1$, $x=-2$ và $x=5$ đều thoả mãn. Vậy 
		\begin{itemchoice}
			\itemch $-1 < 4$.\\
			Mệnh đề sai. 
			\itemch Phương trình $(1)$ có $3$ nghiệm phân biệt.\\
			Mệnh đề đúng.
			\itemch $-1$; $-2$; $5$ đều là các số nguyên.\\
			Mệnh đề đúng.
			\itemch $(-1)^2 + (-2)^2 + 5^2 = 30.$\\
			Mệnh đề sai. 
		\end{itemchoice}	
	}
\end{ex}

\begin{ex}%[0D7V3-4]
	Cho phương trình $\sqrt{3x-2} = 1 + \sqrt{x+7}$. Xác định tính Đúng/Sai của các mệnh đề sau.
	\choiceTF[t]
	{\True Điều kiện xác định của phương trình là $x\ge\dfrac{2}{3}$}
	{Phương trình có hai nghiệm phân biệt}
	{\True Tổng các nghiệm của phương trình là $9$}
	{\True Phương trình có chung tập nghiệm với phương trình $\sqrt{x-5} = x-7$}
	\loigiai{
		Điều kiện xác định $\heva{&3x-2\ge 0\\ &x+7\ge 0} \Leftrightarrow x\ge \dfrac{2}{3}$. \\
		Khi đó phương trình tương đương với 
		\begin{eqnarray*}
			& & 3x - 2  = x+7 + 1 + 2\sqrt{x+7} \Leftrightarrow \sqrt{x+7} = x-5 \\
			&\Leftrightarrow &\heva{&x-5\ge 0\\ &x+7 = (x-5)^2} \Leftrightarrow \heva{&x-5\ge 0\\ &x^2 - 11x + 18 = 0} \Leftrightarrow \heva{&x\ge 5\\ &\hoac{&x=2\\ &x=9}} \Leftrightarrow x=9.
		\end{eqnarray*}
		\begin{itemchoice}
			\itemch Mệnh đề đúng.
			\itemch Mệnh đề sai.
			\itemch Mệnh đề đúng.
			\itemch $\sqrt{x-5} = x-7 \Leftrightarrow \heva{&x-7 \ge 0\\ &x-5 = (x-7)^2} 
			\Leftrightarrow \heva{&x\ge 7\\ &x^2 - 15x + 54 = 0} 
			\Leftrightarrow \heva{&x\ge 7\\ &\hoac{&x=6\\ &x=9}} \Leftrightarrow x = 9$.\\
			Mệnh đề đúng.
		\end{itemchoice}	
	}
\end{ex}


\Closesolutionfile{ans}
\indapan{3}{ans/ans-K10_C6-B4-D1-DS}

\Opensolutionfile{ans}[ans/ans-K10_C6-B4-D1-KQ]
\caukq

\begin{ex}%[0D7C3-6]
	Tam giác $ABC$ vuông tại $A$ có điểm $D$ nằm giữa $A$ và $B$. Biết $BC = 8$, $BD = 3$ và $CD = 6$, tính diện tích tam giác $ABC$ (làm tròn đến chữ số thập phân thứ nhất sau dấu phẩy).
	\shortans{$15{,}7$}
	\loigiai{
		\immini
		{
			Đặt $AD = x$ ($x > 0$). Khi đó theo định lý Py-ta-go ta có
			\begin{itemize}
				\item $AC^2 = CD^2 - AD^2 = 6^2 - x^2 = 36 - x^2$.
				\item $AC^2 = BC^2 - AB^2 = 8^2 - (x+3)^2 = 55 - 6x - x^2$.
			\end{itemize}
		Do đó $36 - x^2 = 55 - 6x - x^2 \Leftrightarrow x = \dfrac{19}{6}$, suy ra 
		\[
			S_{ABC} = \dfrac{1}{2}AB\cdot AC = \dfrac{1}{2}(x+3)\sqrt{36-x^2} \approx 15{,}7. 
		\]
		}
		{
			\begin{tikzpicture}[line join=round, line cap = round, >=stealth, scale=.8,font=\footnotesize,transform shape]
				\foreach \x/\y/\z/\g in
				{
					0/0/A/-135, 3/0/D/-90, 5/0/B/-45, 0/4/C/135
				}
				\draw[fill=black] (\x,\y) circle(1pt) coordinate (\z) ($(\z)+(\g:3.5mm)$) node{$\z$};
				\draw (D)--(C)--(B)--(A)--(C);
			\end{tikzpicture}
		}
	}
\end{ex}

\begin{ex}%[0D7C3-6]
	\immini
	{
		Một chú thỏ ngày nào cũng ra bờ suối ở vị trí $A$ cách cửa hang của mình ở vị trí $B$ là $370$ m để uống nước, sau đó chú thỏ sẽ đến vị trí $C$ cách $A$ $120$ m để ăn cỏ rồi trở về hang. Tuy nhiên, hôm nay sau khi uống nước ở bờ suối, chú thỏ không đến vị trí $C$ như mọi ngày mà chạy đến vị trí $D$ nằm giữa $C$ và $B$ để tìm cà rốt rồi mới trở về hang (xem hình dưới). Biết rằng, tam giác $ABC$ vuông tại $C$, tổng thời gian chú thỏ chạy từ $A$ đến $D$ rồi về hang hết $30$ giây (không kể thời gian tìm cà rốt), tốc độ của chú thỏ trên các đoạn $AD$ và $DB$ lần lượt là $13$ m/s và $15$ m/s. Tính độ dài đoạn $CD$.
	}
	{
		\begin{tikzpicture}[line join=round, line cap = round, >=stealth, scale=.8,font=\footnotesize,transform shape]
			\foreach \x/\y/\z/\g in
			{
				0/0/C/-135, 3/0/D/-90, 5/0/B/-45, 0/4/A/135
			}
			\draw[fill=black] (\x,\y) circle(1pt) coordinate (\z) ($(\z)+(\g:3.5mm)$) node{$\z$};
			\draw (D)--(A)--(B)--(C)--(A);
		\end{tikzpicture}
	}
	\shortans{$50$}
	\loigiai{
		Gọi thời gian chú thỏ chạy trên đoạn $AD$ là $x$ giây ($0<x<30$). Khi đó $AD = 13x$ m.\\ 
		Thời gian chú thỏ chạy trên đoạn $BD$ là $30-x$ giây nên $BD = 15(30-x)$ m.\\
		Theo định lý Py-ta-go ta có $BC = \sqrt{AB^2 - AC^2} = 350$. \\
		Suy ra $CD = BC - BD = 350 - 15(30-x) = 15x - 100$ m.\\
		Áp dụng định lý Py-ta-go cho tam giác $ACD$ ta có 
		\[
			AD^2 = AC^2 + CD^2 \Leftrightarrow (13x)^2 = 120^2 + (15x-100)^2 
			\Leftrightarrow 56x^2 - 3000x + 24400 = 0 \Leftrightarrow \hoac{&x=10\\ &x = \dfrac{307}{5}.} 
		\] 		
		Mà $ 0 < x < 30$ nên $x = 10$, suy ra $CD = 50$.
	}
\end{ex}

\begin{ex}%[0D7C3-4]
	Tính tổng bình phương các nghiệm của phương trình 
	\[
		x^2 - 3x + 10 + 2(x-3)\sqrt{3x+1} = 0.
	\]
	\shortans{$1$}
	\loigiai{
		Ta có phương trình tương đương với
		\begin{eqnarray*}
			& &\Leftrightarrow (x^2 - 6x + 9) + (3x+1) + 2(x-3)\sqrt{3x+1} = 0 \\
			& &\Leftrightarrow (x-3)^2 + \left( \sqrt{3x+1} \right)^2 + 2(x-3)\sqrt{3x+1} = 0\\
			& &\Leftrightarrow \left( x - 3 + \sqrt{3x+1} \right) ^ 2 = 0 \Leftrightarrow x - 3 + \sqrt{3x+1} = 0 \Leftrightarrow \sqrt{3x+1} = 3-x\\
			& & \Leftrightarrow \heva{&3-x\ge 0\\ &3x+1 = (3-x)^2}
			\Leftrightarrow \heva{&x\le 3\\ &x^2 - 9x + 8 = 0}
			\Leftrightarrow \heva{&x\le 3\\ &\hoac{&x=1\\ &x=8}}
			\Leftrightarrow x = 1. 
		\end{eqnarray*}
		Vậy tổng bình phương tất cả các nghiệm của phương trình là $1$.
	}
\end{ex}

\begin{ex}%[0D7C3-2]
	Tính tổng các nghiệm của phương trình $\sqrt{x^2 - 4x - 1} - |2x+1| = 1$ (làm tròn đến một chữ số thập phân sau dấu phẩy).
	\shortans{$-1{,}5$}
	\loigiai{
		Ta xét hai trường hợp:
		\begin{enumerate}[\bf TH1.]
			\item $2x+1\ge 0 \Leftrightarrow x \ge -\dfrac{1}{2}$. Khi đó phương trình trở thành
			\begin{eqnarray*}
				& &\sqrt{x^2 - 4x - 1} - (2x+1) = 1 \Leftrightarrow \sqrt{x^2 - 4x - 1} = 2x+2\\
				& \Leftrightarrow & \heva{&2x+2\ge 0\\ &x^2 - 4x - 1 = (2x+2)^2}
				\Leftrightarrow \heva{&x\ge -1\\ &3x^2 + 12x + 5 = 0}\\
				&\Leftrightarrow &\heva{&x\ge -1\\ &x = \dfrac{-6\pm \sqrt{21}}{3}}
				\Leftrightarrow x = \dfrac{-6+\sqrt{21}}{3} \text{ (thoả mãn)}.
			\end{eqnarray*}
			\item $2x+1< 0 \Leftrightarrow x < -\dfrac{1}{2}$. Khi đó phương trình trở thành
			\begin{eqnarray*}
				& &\sqrt{x^2 - 4x - 1} + (2x+1) = 1 \Leftrightarrow \sqrt{x^2 - 4x - 1} = -2x\\
				& \Leftrightarrow & \heva{&-2x\ge 0\\ &x^2 - 4x - 1 = (-2x)^2}
				\Leftrightarrow \heva{&x\le 0\\ &3x^2 + 4x + 1 = 0}\\
				&\Leftrightarrow &\heva{&x\le 0\\ &\hoac{&x=-1\\ &x=-\dfrac{1}{3}}}
				\Leftrightarrow \hoac{&x=-1\\ &x=-\dfrac{1}{3}.}
			\end{eqnarray*}
			Mà $x< -\dfrac{1}{2}$ nên $x=-1$.
		\end{enumerate}
		Tổng các nghiệm của phương trình là $\dfrac{-6+\sqrt{21}}{3} + (-1) = \dfrac{-9+\sqrt{21}}{3} = \approx -1{,}5$.
	}
\end{ex}

\begin{ex}%[0D7C3-4]
	Tính tổng các nghiệm của phương trình $x^2 -3x + 1 +\sqrt{2x-1} = 0$ (lấy kết quả gồm hai chữ só thập phân sau dấu phẩy).
	\shortans{$1{,}59$}
	\loigiai{
		Điều kiện xác định $x\ge \dfrac{1}{2}$. Khi đó phương trình tương đương 
		\begin{eqnarray*}
			& &\sqrt{2x-1} - x + x^2 - (2x-1) = 0 \Leftrightarrow 
			\left( \sqrt{2x-1} - x \right) + \left( x - \sqrt{2x-1}\right)\left( x + \sqrt{2x-1} \right) = 0 \\
			&\Leftrightarrow & \left( x - \sqrt{2x-1} \right)\left( -1 + x + \sqrt{2x-1} \right)
			= 0 \Leftrightarrow \hoac{&\sqrt{2x-1} = x\\ &\sqrt{2x-1} = 1-x}\\
			&\Leftrightarrow 
			&\hoac{
				&\heva{&x\ge 0\\ & 2x-1 = x^2} &\Leftrightarrow \heva{&x\ge 0\\ &x^2 - 2x + 1 =0} 
				&\Leftrightarrow x = 1\\
				&\heva{&1-x\ge 0\\ & 2x-1 = (1-x)^2} &\Leftrightarrow \heva{&x\le 1\\ &x^2 - 4x + 2 =0} 
				&\Leftrightarrow x = 2 - \sqrt{2}.\\
			}
		\end{eqnarray*}
		Vậy phương trình có tập nghiệm là $S = \{1;2-\sqrt{2}\}$ nên tổng các nghiệm của phương trình là 
		\[
		1 + 2-\sqrt{2 } = 3-\sqrt{2} \approx 1{,}59.
		\]
	}
\end{ex}

\begin{ex}%[0D7C3-4]
	Biết rằng phương trình $(x+1)(x+4) - 5\sqrt{x^2+5x+2} = 6$ có hai nghiệm phân biệt $a$ và $b$ thoả mãn $a < b$. Tính giá trị biểu thức $T = 2a + 12b$.
	\shortans{$10$}
	\loigiai{
		Điều kiện xác định $x^2+5x+2\ge 0$.\\
		Ta có 
		\[
			(x+1)(x+4) - 5\sqrt{x^2+5x+2} = 6 \Leftrightarrow x^2 + 5x - 2 - 3\sqrt{x^2 + 5x+2} = 0.
		\]
		Đặt $t = \sqrt{x^2 + 5x+2}$ ($t\ge 0$) thì $x^2+5x -2 = t^2 - 4$. Khi đó phương trình trở thành 
		\[
			t^2 - 3t - 4 = 0 \Leftrightarrow \hoac{&t = 4 & & \text{ (thoả mãn)}\\ &t = -1 & & \text{ (loại)}.}
		\]
		Với $t=4$ thì
		\[
			\sqrt{x^2+5x+2} = 4\Leftrightarrow x^2 + 5x +2 = 16 \Leftrightarrow x^2 + 5x - 14 = 0 \Leftrightarrow 
			\hoac{&x = 2\\ &x=-7.}
		\]
		Vậy $a=-7$ và $b=2$ nên $T = 2a+14b = 10$.
	}
\end{ex}
\Closesolutionfile{ans}
\indapan{6}{ans/ans-K10_C6-B4-D1-KQ}

\newpage
\subsection{Đề 2}
\Opensolutionfile{ans}[ans/ans-K10_C6-B4-D2]
\caulc

\begin{ex}%[0D7H3-2]
	Nghiệm của phương trình $\sqrt{(x-1)(2x-1)}-2x-1=0$ là
	\choice
	{$x=-\dfrac{7}{2}$}
	{\True $x=0$}
	{$\hoac{&x=0\\&x=-\dfrac{7}{2}}$}
	{$x\in\emptyset$}
	\loigiai{Phương trình đã cho được viết lại $\sqrt{(x-1)(2x-1)}=2x+1$.\\
		Bình phương hai vế của phương trình trên, ta được
		\begin{eqnarray*}
			(x-1)(2x-1)=(2x+1)^2\Rightarrow -2x^2-7x=0\Rightarrow \hoac{&x=0\\ &x=-\dfrac{7}{2}.}
		\end{eqnarray*}
		Thay lần lượt các giá trị trên vào phương trình đã cho, ta thấy $x=0$ thỏa mãn.\\
		Vậy nghiệm của phương trình đã cho là $x=0$.
	}
\end{ex}

\begin{ex}%[0D7H3-2]
	Nghiệm của phương trình $x-\sqrt{3x^2-9x+1}=2$ là
	\choice
	{\True $x=3$}
	{$x=-\dfrac{1}{2}$}
	{$\hoac{&x=-\dfrac{1}{2}\\&x=3}$}
	{$x\in\emptyset$}
	\loigiai{Phương trình đã cho được viết lại $\sqrt{3x^2-9x+1}=x-2$.\\
		Bình phương hai vế của phương trình trên, ta được
		\begin{eqnarray*}
			3x^2-9x+1=(x-2)^2\Rightarrow2x^2-5x-3=0\Rightarrow \hoac{&x=3\\ &x=\dfrac{1}{2}.}.
		\end{eqnarray*}
		Thay lần lượt các giá trị trên vào phương trình đã cho, ta thấy $x=3$ thỏa mãn.\\
		Vậy nghiệm của phương trình đã cho là $x=3$.
	}
\end{ex}



\begin{ex}%[0D7H3-2]
	Tổng các nghiệm của phương trình $\sqrt{x^2+3} = 3x-1$ là
	\choice
	{$0$}
	{\True $1$}
	{$2$}
	{$3$}
	\loigiai{
		Ta có phương trình tương đương với
		\[
		\heva{&3x-1\ge 0\\ &x^2+3 = (3x-1)^2}
		\Leftrightarrow \heva{&x\ge \dfrac{1}{3}\\ &8x^2-6x-2=0}
		\Leftrightarrow \heva{&x\ge\dfrac{1}{3}\\ &\hoac{&x=1\\ &x=-\dfrac{1}{4}}}
		\Leftrightarrow x=1.
		\]	
	}
\end{ex}

\begin{ex}%[0D7H3-2]
	Số nghiệm của phương trình $x - \sqrt{2x^2-3x+1} = 1$ là 
	\choice
	{$0$}
	{$2$}
	{$3$}
	{\True $1$}
	\loigiai{
		Ta có phương trình tương đương với
		\[
		\heva{&x-1\ge 0\\ &2x^2-3x+1 = (x-1)^2}
		\Leftrightarrow \heva{&x\ge 1\\ &x^2-x=0}
		\Leftrightarrow \heva{&x\ge 1\\ &\hoac{&x=0\\ &x=1}}
		\Leftrightarrow x=1.
		\]	
	}
\end{ex}

\begin{ex}%[0D7H3-2]
	Số nghiệm của phương trình $\sqrt{-x^2 + 4x} = 2x-2$ là 
	\choice
	{$0$}
	{$2$}
	{$3$}
	{\True $1$}
	\loigiai{
		Ta có phương trình tương đương với
		\[
		\heva{&2x-2\ge 0\\ &-x^2+4x = (2x-2)^2}
		\Leftrightarrow \heva{&x\ge 1\\ &5x^2 - 12x + 4 = 0}
		\Leftrightarrow \heva{&x\ge 1\\ &\hoac{&x=2\\ &x=\dfrac{2}{5}}}
		\Leftrightarrow x=2.
		\]	
	}
\end{ex}

\begin{ex}%[0D7H3-2]
	Tổng các nghiệm của phương trình $\sqrt{2x-1} = x-2$ là
	\choice
	{$6$}
	{$1$}
	{\True $5$}
	{$2$}
	\loigiai{
		Ta có phương trình tương đương với 
		\[
		\heva{&x-2\ge 0\\ &2x-1 = (x-2)^2}
		\Leftrightarrow \heva{&x\ge 2\\ &x^2 - 6x + 5 = 0}
		\Leftrightarrow \heva{&x\ge 2\\ &\hoac{&x=1\\ &x=5}}
		\Leftrightarrow x=5.
		\]	
	}
\end{ex}

\begin{ex}%[0D7H3-2]
	Số nghiệm nguyên của phương trình $\sqrt{x^2 +10x-5} = 2(x-1)$ là
	\choice
	{\True $0$}
	{$2$}
	{$3$}
	{$1$}
	\loigiai{
		Ta có phương trình tương đương với
		\[
		\heva{&2(x-1)\ge 0\\ &x^2+10x-5 = [2(x-1)]^2}
		\Leftrightarrow \heva{&x\ge 1\\ &3x^2 - 18x  +9 = 0}
		\Leftrightarrow \heva{&x\ge 1\\ &x=3\pm 6}
		\Leftrightarrow x=3+\sqrt{6}.
		\]	
		Vậy phương trình không có nghiệm nguyên.
	}
\end{ex}

\begin{ex}%[0D7C3-6]
	Tam giác $ABC$ vuông tại $A$, có cạnh $AC$ ngắn hơn cạnh $BC$ là $9$ cm. Biết chu vi của tam giác bằng $70$ cm, độ dài đoạn $AC$ là
	\choice
	{$21$ cm}
	{\True $20$ cm}
	{$29$ cm}
	{$28$ cm}
	\loigiai{
		\immini{Gọi độ dài cạnh $AC$ là $x$ cm.\\
			Cạnh $AC$ ngắn hơn cạnh $BC$ là $9$ cm nên độ dài cạnh $BC$ là $(x+9)$ cm.\\
			Ta có $AC^2+AB^2=BC^2$ hay 
			\[
			AB=\sqrt{BC^2-AC^2}=\sqrt{(x+9)^2-x^2}=\sqrt{18x+81}.
			\]
			Chu vi tam giác $ABC$ bằng $70$ cm nên 
			\allowdisplaybreaks
			\begin{eqnarray*}
				AB+BC+CA=70 &\Rightarrow&\sqrt{18x+81}+(x+9)+x=70 \\
				&\Rightarrow&\sqrt{18x+81}=61-2x\\
				&\Rightarrow& 18x+81=61^2-2\cdot 61 \cdot 2x+4x^2\\
				&\Rightarrow& 4x^2-262x+3640=0\\
				&\Rightarrow&\hoac{&x=20\\ &x=\dfrac{91}{2}.}
			\end{eqnarray*}
			Thay lần lượt các giá trị trên vào phương trình đã cho, ta thấy $x=20$ thỏa mãn.\\
			Vậy $AC=20$ cm.
		}{\begin{tikzpicture}[scale=0.7, font=\footnotesize, line join=round, line cap=round, >=stealth]
				\coordinate [label=below left:$A$] (A) at (0,0);
				\coordinate [label=above:$C$](C) at (0,3);
				\coordinate [label=below :$B$] (B) at (2,0);
				\draw (A)--(C)--(B)--(A) ;
				\foreach \diem in {A,B,C} \fill (\diem) circle(1.5pt);
				\draw pic[draw,angle radius=2mm] {right angle = C--A--B};
		\end{tikzpicture}}
	}
\end{ex}

\begin{ex}%[0D7H3-1]
	Nghiệm của phương trình $\sqrt{x^2 - 4x + 4} - \sqrt{(6 - x)(2x - 1)} = 0$ là
	\choice
	{$ x = 5 $}  	
	{\True $ x = \dfrac{2}{3} $}  	
	{$\hoac{& x = \dfrac{2}{3}\\ &x = 5}$}	  
	{ $ x\in \emptyset$}
	\loigiai{
		Điều kiện: $\heva{&x^2 - 4x + 4 \geq 0\\&(6 - x)(2x - 1) \geq 0} \Leftrightarrow \heva{&x  \in \mathbb{R}\\&x \in \left(-\infty; \dfrac{1}{2}\right] \cup  \left[ 6; +\infty\right)} \Leftrightarrow x \in \left(-\infty; \dfrac{1}{2}\right] \cup  \left[ 6; +\infty\right)$.\\
		Ta có 
		\allowdisplaybreaks
		\begin{eqnarray*}
			\sqrt{x^2 - 4x + 4} - \sqrt{(6 - x)(2x - 1)} = 0
			&\Leftrightarrow&\sqrt{x^2 - 4x + 4} = \sqrt{(6 - x)(2x - 1)}\\
			&\Leftrightarrow&x^2 - 4x + 4 =  (6 - x)(2x - 1)\\
			&\Leftrightarrow&x^2 - 4x + 4 =  - 2x^2 + 13x - 6\\
			&\Leftrightarrow&3x^2 - 17x + 10 = 0\\
			&\Leftrightarrow& \hoac{&x = 5 \;\;\text{(loại)}\\&x = \dfrac{2}{3}.}
		\end{eqnarray*}
		Vậy phương trình đã cho có nghiệm là $x = \dfrac{2}{3}$.
	}
\end{ex}

\begin{ex}%[0D7H3-1]
	Nghiệm của phương trình $\sqrt{x^2 - 2x + 4} = \sqrt{2 - x}$ là
	\choice
	{$x = \dfrac{1}{2} $}
	{$x = \dfrac{7}{4} $}	
	{$\hoac{& x = \frac{1}{2} \\& x = \dfrac{7}{4}}$}  
	{\True $x \in \emptyset $}
	\loigiai{
		Điều kiện: $\heva{&x^2 - 2x + 4 \geq 0\\&2 - x \geq 0} \Leftrightarrow \heva{&x  \in \mathbb{R}\\&x \leq 2} \Leftrightarrow x \leq 2$.\\
		Ta có 
		\allowdisplaybreaks
		\begin{eqnarray*}
			\sqrt{x^2 - 2x + 4} = \sqrt{2 - x}
			&\Leftrightarrow&x^2 - 2x + 4 = 2 - x\\
			&\Leftrightarrow&x^2 - x + 2 =  0 \;\;\text{(vô nghiệm)}
		\end{eqnarray*}
		Vậy phương trình đã cho vô nghiệm.
	}
\end{ex}

\begin{ex}%[0D7H3-2]
	Tổng các nghiệm của phương trình $\sqrt{4+2x-x^2} = x-2$ là
	\choice
	{$0$}
	{$1$}
	{$2$}
	{\True $3$}
	\loigiai{
		Bình phương hai vế của phương trình đã cho, ta được
		\begin{eqnarray*}
			4+2x-x^2=(x-2)^2\Rightarrow 4+2x-x^2=x^2-4x+4\Rightarrow 2x^2-6x=0\Rightarrow x=0 \ \text{hoặc} \ x=3.
		\end{eqnarray*}
		Thay lần lượt các giá trị trên vào phương trình đã cho, ta thấy $x=3$ thỏa mãn.\\
		Vậy nghiệm của phương trình đã cho là $x=3$.	
	}
\end{ex}

\begin{ex}%[0D7H3-2]
	Tổng các nghiệm của phương trình $\sqrt{x^2-3x+2}=x-1$ là
	\choice
	{$0$}
	{\True $1$}
	{$2$}
	{$3$}
	\loigiai{
		Bình phương hai vế của phương trình đã cho, ta được
		\begin{eqnarray*}
			x^2-3x+2=(x-1)^2\Rightarrow x^2-3x+2=x^2-2x+1\Rightarrow -x+1=0\Rightarrow x=1.
		\end{eqnarray*}
		Thay $x=1$ vào phương trình đã cho, ta thấy $x=1$ thỏa mãn.\\
		Vậy nghiệm của phương trình đã cho là $x=1$.	
	}
\end{ex}

\Closesolutionfile{ans}
\indapan{6}{ans/ans-K10_C6-B4-D2}
\cauds
\Opensolutionfile{ans}[ans/ans-K10_C6-B4-D2-DS]

\begin{ex}%[0D7V3-2]
	Cho phương trình $\sqrt{2x^2 + x - 6} = x+2\quad (1)$. Xác định tính Đúng/Sai của các mệnh đề sau.
	\choiceTF[t]
	{Bình phương hai vế của phương trình $(1)$ ta được phương trình $x^2 + 3x - 10 = 0$}
	{Điều kiện xác định của phương trình là $x\ge 3$}
	{\True Phương trình $(1)$ có hai nghiệm phân biệt}
	{\True Tổng các nghiệm của phương trình $(1)$ là $3$}
	\loigiai{
		Bình phương hai vế của phương trình $(1)$ ta được 
		\[
		2x^2 + x - 6  = (x+2)^2 \Leftrightarrow x^2 - 3x - 10 = 0 \Leftrightarrow \hoac{&x=5\\ & x = -2.}
		\]
		Thử lại ta thấy $x=5$ và $x=-2$ đều thoả mãn. Vậy
		\begin{itemchoice}
			\itemch Mệnh đề sai. 
			\itemch Phương trình xác định khi và chỉ khi
			\[
				2x^2 + x - 6 \ge 0 \Leftrightarrow (x+2)(x-3)\ge 0 \Leftrightarrow \hoac{&x\ge 3\\ &x\le -2.}
			\] 
			Mệnh đề sai.
			\itemch Mệnh đề đúng. 
			\itemch $5 + (-2) = 3$.\\
			Mệnh đề đúng. 
		\end{itemchoice}	
	}
\end{ex}

\begin{ex}%[0D7V3-1]
	Cho phương trình $\sqrt{x^2 - 4x - 5} = \sqrt{2x^2 + 3x + 1} \quad (1)$. Xác định tính Đúng/Sai của các mệnh đề sau.
	\choiceTF[t]
	{Bình phương hai vế của phương trình $(1)$ ta được phương trình $x^2 - 7x + 6 = 0$}
	{$x=-1$ là nghiệm của phương trình $(1)$}
	{Tổng các nghiệm của phương trình $(1)$ là $-5$}
	{Phương trình sao có nhiều hơn $1$ nghiệm}
	\loigiai{
		Bình phương hai vế của phương trình $(1)$ ta được 
		\[
		x^2 - 4x - 5 = 2x^2 + 3x + 1 \Leftrightarrow x^2+7x+6 = 0 \Leftrightarrow \hoac{&x=-1\\ &x=-6.}
		\]
		Thử lại ta thấy $x=-1$ và $x=-6$ đều thoả mãn. Vậy
		\begin{itemchoice}
			\itemch Mệnh đề sai. 
			\itemch Mệnh đề đúng. 
			\itemch Tổng các nghiệm của phương trình $(1)$ là $-7$.
			\\ Mệnh đề sai.
			\itemch Phương trình $(1)$ có $2$ nghiệm. 
			\\ Mệnh đề đúng.
		\end{itemchoice}	
	}
\end{ex}

\begin{ex}%[0D7V3-2]
	Cho hai phương trình $\sqrt{5x+10} = 8 - x \quad (1)$ và $\sqrt{3x^2 - 9x + 1} = x - 2 \quad (2)$. Xác định tính Đúng/Sai của các mệnh đề sau.
	\choiceTF[t]
	{\True Phương trình $(1)$ có duy nhất một nghiệm}
	{Tổng số nghiệm của hai phương trình là $1$}
	{\True Hai phương trình có chung tập nghiệm}
	{Tổng tất cả các nghiệm của hai phương trình là $6$}
	\loigiai{
		Ta có 
		\begin{itemize}
			\item $(1) \Leftrightarrow \heva{&8-x \ge 0\\ &5x+10 = (8-x)^2} 
			\Leftrightarrow \heva{&x\le 8\\ &x^2 - 21x + 54 = 0} \Leftrightarrow \heva{& x\le 8\\ & \hoac{&x = 3\\ &x = 18}}
			\Leftrightarrow x=3$.
			\item $(2) \Leftrightarrow \heva{&x-2\ge 0\\ & 3x^2 - 9x + 1 = (x-2)^2} \Leftrightarrow 
			\heva{&x\ge 2\\ &2x^2 - 5x - 3 = 0} \Leftrightarrow \heva{&x\ge 2\\ &\hoac{&x=3 \\ & x = -\dfrac{1}{2}}}
			\Leftrightarrow x = 3$.
		\end{itemize}
		Vậy
		\begin{itemchoice}
			\itemch Mệnh đề đúng. 
			\itemch Mỗi phương trình đều có duy nhất nghiệm nên tổng số nghiệm của hai phương trình là $2$.\\
			Mệnh đề sai.
			\itemch Hai phương trình đều có tập nghiệm là $S = \{1\}$.\\
			Mệnh đều đúng.
			\itemch Tổng tất cả các nghiệm của hai phương trình là $1 + 1 = 2$.\\
			Mệnh đề sai.
		\end{itemchoice}	
	}
\end{ex}

\begin{ex}%[0D7V3-5]
	Cho phương trình $(x-2) \sqrt{2x^2+4} = x^2 - 4$. Xác định tính Đúng/Sai của các mệnh đề sau.
	\choiceTF[t]
	{\True Số nghiệm phân biệt của phương trình là một số lẻ}
	{\True Tất cả các nghiệm của phương trình đều là số tự nhiên chẵn}
	{Tổng tất cả các nghiệm của phương trình là $10$}
	{Không tồn tại cấp số cộng chứa tất cả các nghiệm của phương trình}
	\loigiai{
		Ta có phương trình đã cho tương đương với 
		\begin{eqnarray*}
			& &\Leftrightarrow (x-2)\sqrt{2x^2 + 4} = (x-2)(x+2) 
			\Leftrightarrow \hoac{&x - 2 = 0\\ &\sqrt{2x^2 +4} = x+2}\\
			& &\Leftrightarrow \hoac{&x=2\\ & \heva{&x+2\ge 0\\ &2x^2 + 4 = (x+2)^2}} 
			\Leftrightarrow \hoac{&x=2\\ & \heva{&x\ge -2\\ &x^2 - 4x = 0}}
			\Leftrightarrow \hoac{&x=2\\ &\heva{&x\ge -2\\ &\hoac{&x=0\\ &x=4}}}
			\Leftrightarrow \hoac{&x=0 \\ &x=2\\ &x=4.}
		\end{eqnarray*}
		Vậy
		\begin{itemchoice}
			\itemch Phương trình có $3$ nghiệm phân biệt và $3$ là số lẻ.\\
			Mệnh đề đúng.
			\itemch $0$, $2$, $4$ đều là số tự nhiên chẵn.\\
			Mệnh đề đúng.
			\itemch $0+2+4 = 6$.\\
			Mệnh đề sai.
			\itemch Do $0+4=2\cdot 2$ nên $0$, $2$, $4$ lập thành một cấp số cộng.\\
			Mệnh đề sai.
		\end{itemchoice}	
	}
\end{ex}

\Closesolutionfile{ans}
\indapan{3}{ans/ans-K10_C6-B4-D2-DS}

\Opensolutionfile{ans}[ans/ans-K10_C6-B4-D2-KQ]
\caukq

\begin{ex}%[0D7C3-6]
	\immini
	{
		Ba thành phố $A$, $B$, $C$ nằm tại ba vị trí lập thành một tam giác vuông tại $B$. Khoảng cách $AB$ là $6$ km, khoảng cách $BC$ là $15$ km. Để nối gần hai thành phố $C$ và $A$, trên đường đi từ $C$ đến $B$ người ta xây xựng một bến tàu tại $M$ cho mọi người có thể đi tàu theo đường sông tới $A$. Bác An di chuyển theo tuyến đường mới, đi bộ từ $C$ đến $M$ với vận tốc $10$ km/h và đi tàu từ $M$ tới $A$ với vận tốc $30$ km/h thì hết $1$ giờ $16$ phút tới được $A$. Tính tổng quãng đường bác An đi từ $A$ tới $C$ theo đơn vị ki-lô-mét.
	}
	{
		\begin{tikzpicture}[line join=round, line cap = round, >=stealth, scale=.8,font=\footnotesize,transform shape]
			\foreach \x/\y/\z/\g in
			{
				0/0/B/-135, 3/0/M/-90, 5/0/C/-45,0/2/A/90
			}
			\draw[fill=black] (\x,\y) circle(1pt) coordinate (\z) ($(\z)+(\g:3.5mm)$) node{$\z$};
			\draw (M)--(A)--(B)--(C);
		\end{tikzpicture}
	}
	\shortans{$17$}
	\loigiai{
		Đổi $1$ giờ $6$ phút $=$ $\dfrac{31}{30}$ giờ.\\
		Gọi thời gian đi quãng đường $CM$ của bác An là $x$ (giờ, $0 < x < \dfrac{31}{30}$) thì $CM = 10x$ hay $BM = 15 -10x$.\\
		Thời gian đi quãng đường $MA$ của bác An là $\dfrac{31}{30} - x$ nên $AM = 30\left( \dfrac{31}{30} - x \right) = 31 - 30x$.\\
		Theo định lý Py-ta-go ta có $AM^2 = AB^2 + BM^2$ hay 
		\[
		(31 - 30x)^2 = 6^2 + (15-10x)^2 \Leftrightarrow 800x^2 - 1560x + 700 = 0 \Leftrightarrow  \hoac{&x = 0{,}7 & &\text{ (thoả mãn)}\\ &x=1{,}25 & & \text{ (loại)}.}
		\]
		Do đó tổng quãng đường bác An đi được là $CM + AM = 31 - 20x = 17$ km.
	}
\end{ex}

\begin{ex}%[0D7C3-6]
	Ba thành phố $A$, $T$, $B$ cùng với các con đường nối các thành phố tạo thành một tam giác vuông tại $T$, khảng cách $AB$ là $100$ km. Tại cùng một thời điểm, ô tô thứ nhất đi từ $A$ tới $T$ với vận tốc $55$ km/h và ô tô thứ hai đi từ $B$ tới $T$ với vận tốc $45$ km/h. Tại thời điểm xe thứ nhất cách $T$ $14$ km thì xe thứ hai cách $T$ $6$ km, hỏi lúc này hai xe đã di chuyển bao nhiêu giờ? 
	\shortans{$1{,}2$}
	\loigiai{
		\immini
		{
			Gọi thời gian hai xe chạy là $t$ (giờ, $t>0$). \\
			Khi đó $AT = 55t + 14$ km và $BT = 45t + 8$ km.\\
			Do tam giác $ABT$ vuông tại $T$ và $AB = 100$ km nên
			\begin{eqnarray*}
				& &AB^2 = AT^2 + BT^2\Leftrightarrow (55t+14)^2 + (45t+6)^2 = 100^2 \\
				&\Leftrightarrow &5050t^2 + 2080t - 9768 = 0 \Leftrightarrow \heva{&t = \dfrac{6}{5}\\ & t = - \dfrac{814}{505}.}
			\end{eqnarray*}
			Do $t>0$ nên $t = \dfrac{6}{5} = 1{,}2$.
		}
		{
			\begin{tikzpicture}[line join=round, line cap = round, >=stealth, scale=.8,font=\footnotesize,transform shape]
				\foreach \x/\y/\z/\g in
				{
					0/0/A/-135, 4/0/T/45, 4/-3/B/-90
				}
				\draw[fill=black] (\x,\y) circle(1pt) coordinate (\z) ($(\z)+(\g:3.5mm)$) node{$\z$};
				\draw (A)--(T)--(B)--(A);
			\end{tikzpicture}
		}
	}
\end{ex}

\begin{ex}%[0D7C3-3]
	Tìm số giá trị nguyên của tham số $m$ thuộc đoạn $[0;2024]$ thoả mãn phương trình $\sqrt{2x^2 - 2x - 2m} = x-2$ có nghiệm.
	\shortans{$2023$}
	\loigiai{
		Ta có 
		\[
		\sqrt{2x^2 - 2x - 2m} = x-2 \Leftrightarrow \heva{&x-2\ge 0\\ &2x^2 - 2x - 2m = (x-2)^2}
		\Leftrightarrow \heva{&x\ge 2\\ &x^2+2x-4 = 2m.}
		\]
		Do đó phương trình $\sqrt{2x^2 - 2x - 2m} = x-2$ có nghiệm khi và chỉ khi phương trình $x^2+2x-4 = 2m$ có nghiệm $x_0\ge 2$, điều này tương đương với 
		\[
		2m \ge \min\limits_{x\ge 2} \left(x^2+2x-4\right)
		\]
		Xét tam thức bậc hai $f(x) = x^2+2x-4$ trên $\left[2;+\infty \right)$ có $a = 1>0$ và hoành độ đỉnh $x = -\dfrac{b}{2a} =-1 < 2$ nên ta có 
		\[
		f(x) \ge f(2) = 4,\, \forall x\ge 2.
		\]
		Suy ra $\min\limits_{x\ge 2} \left(x^2+2x-4\right) = 4$ hay $2m\ge 4 \Leftrightarrow m\ge 2$.\\
		Kết hợp với $m\in[0;2024]$ ta thấy có $2023$ giá trị nguyên của $m$ thoả mãn.
	}
\end{ex}

\begin{ex}%[0D7C3-6]
	\immini{
		Một đài quan sát $O$ cách ba vị trí $A$, $B$, $C$ như hình vẽ dưới đây thoả mãn $OB = x$ km, $OC = x+1$ km và $OA = 2$ km. Tìm $x$ biết khoảng cách từ vị trí $A$ đến vị trí $C$ gấp đôi khoảng cách từ vị trí $A$ đến vị trí $B$ và khoảng cách từ $O$ đến $B$ ngắn hơn khoảng cách từ $O$ đến $A$.
	}{
		\begin{tikzpicture}[line join = round, line cap = round, scale=.8] 
			\coordinate[label = left:$B$] (B) at (0,0); 
			\coordinate[label = right:$C$] (C) at ($(B)+(6,0)$); 
			\coordinate[label = above:$A$] (A) at ($(B)+(0,4)$);
			\coordinate[label = below:$O$] (O) at ($(B)+(2,0)$);
			\draw[fill = gray!50] ($(B)!5pt!(C)$)--($(B)!5pt!(C)+(B)!5pt!(A)-(B)$)--($(B)!5pt!(A)$)--(B)--cycle;
			\draw (B)--(C)--(A)--cycle (A)--(O); 
			\foreach \x in {B,C,A,O} \fill[black] (\x) circle (1.5pt);
			\draw pic[draw,black,angle radius=3mm] {angle = C--O--A} ($($(O)!18mm!(C)$)!.6!($(O)!13mm!(A)$)$)node[rotate=0]{$120^\circ$};
		\end{tikzpicture}
	}
	\shortans{$1$}
	\loigiai{
		Áp dụng định lý cosin trong tam giác $AOC$ ta được 
		\[
			AC^2=AO^2+OC^2-2\cdot AO\cdot OC \cdot\cos\widehat{AOC}=4+(x+1)^2-4(x+1)\cdot\cos120^\circ=x^2+4x+7.
		\]
		Suy ra $AC=\sqrt{x^2+4x+7}$.\\
		Áp dụng định lý cosin trong tam giác $AOB$ ta được 
		\[
			AB^2=AO^2+OB^2-2\cdot AO\cdot OB\cos\widehat{AOB}=4+x^2-4x\cos60^\circ=x^2-2x+4.
		\]
		Suy ra $AB=\sqrt{x^2-2x+4}$.\\
		Theo giả thiết ta có $AC=2BO \Leftrightarrow 2\sqrt{x^2-2x+4}=\sqrt{x^2+4x+7}.\quad (*)$\\
		Giải phương trình $2\sqrt{x^2-2x+4}=\sqrt{x^2+4x+7}.\quad(*)$\\
		Bình phương hai vế của phương trình $(*)$ ta được phương trình $4(x^2-2x+4)=x^2+4x+7$. \\
		Ta có $4(x^2-2x+4)=x^2+4x+7 \Leftrightarrow 3x^2-12x+9=0 \Leftrightarrow 
		\hoac{&	x=1\\&x=3.}$\\
		Thay hai giá trị $x=1$ và $x=3$ vào phương trình $(*)$ ta thấy cả hai giá trị đều thõa  mãn, nhưng $OB<OA$ nên $x=1$.		
	}
\end{ex}

\begin{ex}%[0D7C3-4]
	Tìm số nghiệm của phương trình $\sqrt{x + 2\sqrt{x-1}} = x+\dfrac{1}{4}$.
	\shortans{$1$}
	\loigiai{
		Ta có $x+2\sqrt{x-1} = x-1 + 2\sqrt{x-1} + 1 = \left( \sqrt{x-1} + 1 \right)^2$ và $\sqrt{x-1} + 1 > 0$ với mọi $x\ge 1$ nên phương trình đã cho tương đương với
		\begin{eqnarray*}
			& &\sqrt{x-1} + 1 = x + \dfrac{1}{4} 
			\Leftrightarrow 4\sqrt{x-1} = 4x-3 \\
			&\Leftrightarrow  &\heva{&4x-3\ge 0\\ &16(x-1) = (4x-3)^2}
			\Leftrightarrow \heva{&x\ge \dfrac{3}{4}\\ &16x^2 - 40x + 25 = 0}\\
			&\Leftrightarrow & \heva{&x\ge \dfrac{3}{4}\\ & x = \dfrac{5}{4}} \Leftrightarrow x = \dfrac{5}{4}.
		\end{eqnarray*}
		Vậy phương trình có nghiệm duy nhất.
	}
\end{ex}

\begin{ex}%[0D7C3-4]
	Tính tổng bình phương các nghiệm của phương trình 
	\[
	\sqrt{x^2 - 3x + 3} + \sqrt{x^2 - 3x + 6} = 3.
	\]
	\shortans{$5$}
	\loigiai{
		Đặt $t = \sqrt{x^2 - 3x + 3}$ ($t\ge 0$) thì $x^2 - 3x + 6 = t^2 + 3$. Khi đó phương trình trở thành 
		\[
		t + \sqrt{t^2 + 3} = 3
		\Leftrightarrow \heva{&3-t\ge0 \\ &t^2 + 3 = (3-t)^2}
		\Leftrightarrow \heva{&t\le 3\\ & t = 1} 
		\Leftrightarrow t = 1.
		\]
		Với $t = 1$ thì 
		\[
		\sqrt{x^2 - 3x + 3} = 1 \Leftrightarrow x^2 - 3x + 3 = 1 \Leftrightarrow x^2 - 3x + 2 = 0 \Leftrightarrow \hoac{
			&x=1\\ &x=2.}
		\]
		Vậy tổng bình phương các nghiệm của phương trình là $1^2 + 2^2 = 5$.
	}
\end{ex}

\Closesolutionfile{ans}
\indapan{6}{ans/ans-K10_C6-B4-D2-KQ}